\documentclass[11pt,letterpaper]{article}
\usepackage[T1]{fontenc}
\usepackage{lmodern,amsmath,amssymb,amsthm,mathtools,booktabs,longtable,microtype}
\usepackage[margin=.94in,headheight=14pt]{geometry}
\usepackage{xcolor,enumitem,fancyhdr}
\usepackage[colorlinks=true,linkcolor=teal,citecolor=teal,urlcolor=teal]{hyperref}
\definecolor{ink}{HTML}{18324A}
\pagestyle{fancy}\fancyhf{}
\fancyhead[L]{\small\sffamily\color{ink} Exact local fewnomial bounds}
\fancyhead[R]{\small 27 November 2026}
\fancyfoot[C]{\thepage}
\setlength{\emergencystretch}{2em}
\newtheorem{theorem}{Theorem}[section]
\newtheorem{lemma}[theorem]{Lemma}
\newtheorem{proposition}[theorem]{Proposition}
\newtheorem{corollary}[theorem]{Corollary}
\theoremstyle{definition}\newtheorem{definition}[theorem]{Definition}
\theoremstyle{remark}\newtheorem{remark}[theorem]{Remark}
\newcommand{\Q}{\mathbb Q}
\newcommand{\Z}{\mathbb Z}
\newcommand{\Qp}{\mathbb Q_p}
\newcommand{\Cp}{\mathbb C_p}
\newcommand{\Fp}{\mathbb F_p}
\newcommand{\vp}{v_p}
\newcommand{\ds}{s_p}
\newcommand{\ord}{\operatorname{ord}}
\newcommand{\supp}{\operatorname{supp}}
\newcommand{\remn}{\operatorname{rem}}
\newcommand{\repo}{https://github.com/VladimirReshetnikov/ProveIt/blob/b82291c1f75389938c7464ef161907d7d5e75b2a/Analysis/GeneralizedRiemannHypothesis}
\title{\color{ink}Finite residue classification of sparse polynomial roots\\
\large Exact local bounds, a sharp odd-prime staircase, and binary disk profiles}
\author{Research supplement to the ProveIt GRH transfer program}
\date{27 November 2026}
\begin{document}
\maketitle
\begin{abstract}
For a prescribed finite set of integer exponents, we determine the largest
possible number of geometric roots of a polynomial in a positive-depth
$p$-adic disk about one. A finite matrix gives the exact answer, including
the endpoints at which the answer changes. More strongly, above any fixed
positive depth the entire optimal profile depends only on finitely many
residue digits of the exponents, with an explicit precision independent of
their sizes. Thus the sharp local fewnomial constant, optimized over both
supports and coefficients, is computable for every prime, term count and
positive rational depth. The initial uniform cutoff comes from Lenstra's
local theorem; the exact classification comes from minimum-valuation bases
of a quotient-algebra matrix. Inside the exponential disk we give a
depth-independent cutoff and residue precision by Frobenius compression.
We also prove the sharp formula
$p^2+\lfloor(t-p-2)/2\rfloor$ for $p+2\le t\le2p-1$, and exact binary
profiles at five, six and seven terms. Rational coefficients attain every
optimum; for at least two terms, full support, irreducibility and simple
roots can be imposed.
These results extend the repository's local transfer machinery and do not
assume RH or GRH. A separate companion paper proves a positive-tolerance
growth law for finite torsion packets.
\end{abstract}
{\small\tableofcontents}

\section{Problem, scope, and contribution}
Fix a prime $p$ and normalize $\vp(p)=1$. Roots are counted in the
algebraically closed complete field $\Cp$, with multiplicity unless
otherwise stated. For $r>0$ set
\[
 D_r=\{x\in\Cp:\vp(x-1)\ge r\}.
\]
For a set $E=\{e_1,\ldots,e_t\}$ of distinct nonnegative integers, define
\begin{align*}
 C_E(r)&=\max_{(c_1,\ldots,c_t)\ne0}
 \#\{x\in D_r:\textstyle\sum_i c_i x^{e_i}=0\},\\
 B_p(t,r)&=\max_{|E|\le t}C_E(r).
\end{align*}
The coefficients initially range over $\Cp$. The theorems below show
that rational coefficients suffice and that each maximum is attained.
Multiplying a polynomial by a monomial does not change its roots in
$D_r$, so the exponents may be translated to make their minimum zero.

The integer exponents may be extremely large. A method whose size grows
with $\max E$ therefore misses the sparse nature of the problem. Our
matrix has a number of columns bounded in terms of $p,t,r$, and not in
terms of $\max E$. The strongest structural statement is the following.

\begin{theorem}[Finite residue classification]\label{thm:headline}
Fix $p,t$ and a positive rational $r_0$. There is an explicit integer
$P=P(p,t,r_0)\ge1$ with these properties.
\begin{enumerate}[label=\textup{(\alph*)}]
\item If two $t$-element supports have the same residue multiset modulo
$p^P$, then $C_E(r)=C_{E'}(r)$ for every real $r\ge r_0$.
\item Every $C_E(r)$ is given by a finite exact determinant calculation.
The profile is a nonincreasing step function with finitely many rational
transition depths on $[r_0,\infty)$; the larger count is retained at a
transition depth.
\item The sharp constant $B_p(t,r)$ is computable from $p,t$ and positive
rational $r$. It is attained by an integer polynomial with exactly $t$
nonzero terms, and for $t\ge2$ the polynomial can be required to be
irreducible after removal of a monomial factor.
\end{enumerate}
An explicit choice of $P$ is given in Theorem~\ref{thm:residues-all}.
\end{theorem}

Here ``computable'' asserts a terminating exact algorithm. The finite
enumeration over all residue multisets can be enormous. No claim of
practical efficiency in all three parameters is intended. For a fixed
$p,t,r$, the prescribed-support computation uses a fixed-size matrix
whose entries can be formed from exponents in binary.

\paragraph{Relation to the repository.}
The snapshot used throughout is
\href{https://github.com/VladimirReshetnikov/ProveIt/commit/b82291c1f75389938c7464ef161907d7d5e75b2a}{commit \texttt{b82291c1f753}};
the full identifier and source hashes are in the archive manifest.
The article's singular-modulus section develops the moment recurrence
and local root bounds; subsequent reports give a residue-rank rule for
at most $p$ exponents and an open-disk formula through $p+2$ terms
\cite{ProveIt,RankPacket,SharpPacket}.
The present work removes the term-count restriction from the exact
algorithm, treats every positive disk depth, and gives new sharp
formulas beyond that earlier range. The ambient project's RH input is
the Dirichlet statement in its pinned source, whose checks are reported
separately \cite{Sernau,Sutherland}. None of the proofs here calls that
input theorem.

\paragraph{Relation to classical work.}
Lenstra proves a uniform bound in exactly the geometric local setting
used here \cite[Theorem 3 and Section 3]{Lenstra}. We use that bound as a
cutoff, then optimize exactly within it. The theory of exponential
polynomials and root counts over local fields has a much broader history
\cite{PoortenRumely,AvendanoKrick}. Minimum-valuation bases are a standard
linear-algebra manifestation of valuated matroids \cite{Battistella}.
Our proof of the needed basis fact is included, so no tropical machinery
is required. The claimed new content is the explicit finite residue
classification, its application to sharp local constants, and the
formulas proved below. We have not established publication priority.

\section{A uniform finite cutoff at every positive depth}
We first state precisely the classical input needed for the all-depth
result. Put $k=t-1$ and define
\[
 a_k(m)=\sum_{a\ge1}\min\{k,\lfloor m/p^a\rfloor\},\qquad
 W_k(r)=\max_{0\le i\le k}\{ri-\vp(i!)\}.
\]
Only finitely many summands in $a_k(m)$ are nonzero. It is the largest
valuation of a product of at most $k$ distinct integers from $1$ to $m$.
Set
\begin{equation}\label{eq:lenstra-cutoff}
 N_L(p,t,r)=\max\{m\ge0:rm-a_k(m)\le W_k(r)\}.
\end{equation}
For $t=1$ put $N_L=0$.

\begin{lemma}[Lenstra's local cutoff]\label{lem:lenstra}
Every nonzero polynomial with at most $t$ terms has at most
$N_L(p,t,r)$ roots in $D_r$, counted geometrically with multiplicity.
The bound is independent of the coefficient field and of the exponents.
\end{lemma}

This is the explicit cutoff in \cite[Section 3]{Lenstra}, expressed using
the valuation of his integer $d_k(m)$. Its coefficient argument uses
integer-valued interpolation: for any $t$ integer nodes and $m\ge t$, the values of
$\binom{T}{m}$ are a rational linear combination of those of
$\binom{T}{i}$, $0\le i<t$, with the $i$th denominator dividing
$d_k(m)/i!$. Applying this relation to the coefficients of $f(1+Y)$ at
their last Newton minimum gives the inequality defining
\eqref{eq:lenstra-cutoff}.

The cutoff is effectively finite. Indeed
$a_k(m)\le k\lfloor\log_p m\rfloor$. Choose $J$ such that
\[
 rp^J-kJ>W_k(r),\qquad rp^J(p-1)>k.
\]
These inequalities persist for larger $J$. Every $m\ge p^J$ is then
excluded from \eqref{eq:lenstra-cutoff}, by considering its interval
$p^j\le m<p^{j+1}$. Checking the finitely many smaller integers computes
$N_L$. Also $N_L\ge t-1$, because $a_{t-1}(t-1)=\vp((t-1)!)$.

\section{The exact minimum-basis rule}
Write
\[
 F_E(T)=\prod_{e\in E}(T-e),\qquad
 A_E=\Q[T]/F_E(T).
\]
Use $1,T,\ldots,T^{t-1}$ as its basis. Let $q_k\in\Z^t$ be the
coefficient vector of the remainder of the falling factorial
\[
 (T)_k=T(T-1)\cdots(T-k+1),\qquad (T)_0=1,
\]
on division by the monic polynomial $F_E$. All these vectors are
integral. The matrix with columns $q_0,\ldots,q_{t-1}$ is unit
triangular, so its determinant is one.

For $N\ge t-1$ and a $t$-element subset $J\subseteq\{0,\ldots,N\}$,
let $q_J$ be its column submatrix and set
\begin{equation}\label{eq:basis-cost}
 \nu_J(r)=\vp(\det q_J)-\sum_{j\in J}\vp(j!)+r\sum_{j\in J}j,
\end{equation}
with valuation $+\infty$ for a zero determinant. Let
$\mathcal B(r)$ be the nonempty collection of subsets minimizing this
quantity. The initial subset $J_0=\{0,\ldots,t-1\}$ has finite cost.

\begin{theorem}[Exact local formula]\label{thm:exact}
Fix $r>0$ and choose any uniform root cutoff $N$ for $t$-term
polynomials on $D_r$, for example $N=N_L(p,t,r)$. Then
\begin{equation}\label{eq:exact}
             C_E(r)=\max\bigcup_{J\in\mathcal B(r)}J.
\end{equation}
All quantities are rational numbers and integer valuations when $r$ is
rational. A dual row of a minimizing basis constructs rational
coefficients attaining the maximum.
\end{theorem}

\begin{proof}
For rational $r$, choose $b\in\overline\Q_p$ with $\vp(b)=r$.
Every coefficient vector of $f(X)=\sum_{e\in E}c_eX^e$ gives the
functional
$L(H)=\sum_e c_eH(e)$ on $A_E$; conversely every functional is
obtained, by the invertibility of the Vandermonde matrix. The
coefficient of $Y^k$ in $f(1+bY)$ is
\[
                   L(q_k)b^k/k!.
\]
The number of roots in $D_r$ is its largest coefficient index of
minimum valuation. By the cutoff, that index is at most $N$ for every
functional. Thus the columns
$V_k=q_kb^k/k!$, $0\le k\le N$, contain all the needed information.

Choose a minimum-determinant basis $B$ of this weighted matrix, and
write $U=B^{-1}V$. Every entry of $U$ is integral: replacing one basis
column by a column with a negative-valuation coordinate would strictly
decrease the determinant valuation. A column of $U$ has a unit entry
if and only if its index belongs to some minimum basis. One direction
is the corresponding one-column exchange. For the other, if an
alternative minimum determinant is expanded in the $B$ coordinates,
at least one product in the determinant expansion consists entirely
of units; each of its columns consequently contains a unit.

For an arbitrary functional, put $a=L B$. The minimum valuation of
its values on the columns of $V$ equals $\min_i\vp(a_i)$, since the
basis columns supply exactly the entries of $a$ and all coordinates
of $U$ are integral. A column with no unit coordinate cannot attain
this minimum. This proves the upper bound in \eqref{eq:exact}.

Let $j$ be the largest index belonging to a minimum basis, and choose
one containing it. Take its dual row at $j$. Its value at column $j$
is a unit, all values are integral, and every later value has positive
valuation. Hence its last minimum is exactly $j$. The weighted inverse
differs from the inverse of the rational matrix $q_J$ by diagonal
scalars. Omitting the common scalar in this row gives a rational
functional $L$, and rational Vandermonde inversion gives rational
$c_e$. The uniform cutoff ensures that omitted columns beyond $N$
cannot create a new last minimum.

For irrational $r$, first use the full Taylor matrix through degree
$\max E$. Its affine basis costs have rational crossing points.
Rational depths just below and just above $r$ in the same chamber give
the same optimum; the inclusions between their disks prove the formula
at $r$ too. Any minimum basis using an index greater than $N$ would
provide a witness with that many roots at a slightly greater rational
depth, hence in $D_r$, contradicting the assumed cutoff. Thus all full
minimum bases lie in the truncated matrix, proving the asserted
formula with $N$ as well.
\end{proof}

\begin{proposition}[Full support and irreducible witnesses]\label{prop:witness}
Every maximum in Theorem~\ref{thm:exact} is attained by a rational
polynomial in which all exponents of $E$ occur. If $|E|\ge2$, after
factoring out $X^{\min E}$, the polynomial may be chosen integral and
Eisenstein at an auxiliary prime different from $p$. All its nonzero
roots are then simple. For $|E|=1$ there are no roots in the disk.
The same assertions hold for the open exponential disk below.
\end{proposition}

\begin{proof}
For $|E|=1$ take any nonzero monomial; henceforth assume $|E|\ge2$.
Clear the denominators of a rational witness. A sufficiently small
$p$-adic perturbation of its finitely many allowed coefficients
preserves the last minimum at the chosen depth: the coefficients of
each $(1+bY)^e$ have nonnegative valuation, so a perturbation divisible
by a sufficiently high power of $p$ cannot change that minimum or its
last index. For an irrational depth choose a nearby rational depth in
the same chamber. For an open disk choose the final chamber next to
its boundary, as below.

Remove the monomial factor. Choose an auxiliary prime $\ell\ne p$.
The Chinese remainder theorem allows every allowed coefficient to
stay in its prescribed neighborhood modulo a high power of $p$,
while the leading coefficient is a unit modulo $\ell$, the constant
is $\ell$ modulo $\ell^2$, and every other coefficient is zero modulo
$\ell$. Unsupported coefficients already vanish. Each supported
coefficient can be required to be nonzero. The resulting polynomial
is Eisenstein at $\ell$, hence irreducible and separable over $\Q$.
The monomial case is immediate.
\end{proof}

\subsection{A direct exact algorithm}
One need not enumerate all minors to find a minimum basis. Start with
$J_0$. In current unweighted basis coordinates write $u_{ij}$ for
column $j$, and let the $i$th basis index be $b_i$. Replacing $b_i$ by
$j$ changes the objective by
\[
 \vp(u_{ij})+r(j-b_i)-\vp(j!)+\vp(b_i!).
\]
Perform a rational pivot whenever this is negative. The objective
strictly decreases and there are finitely many bases. When no such
exchange exists, the basis is globally minimal: each term of any
alternative determinant expansion has nonnegative relative cost.
The indices occurring in minimum bases are the current basis indices
and the columns admitting a zero-cost exchange. This also gives the
dual row needed for a witness. The supplied program implements these
operations with exact fractions. A parametric version tracks the
affine exchange inequalities and hence every transition depth.

\section{Only finitely many exponent digits matter}
\begin{theorem}[All positive depths above a prescribed floor]\label{thm:residues-all}
Fix $r_0>0$ rational, let $N=N_L(p,t,r_0)$, and put
\begin{align}
 D_0&=\sum_{i=0}^{t-1}\vp(i!),&
 D_{\max}&=\sum_{i=N-t+1}^{N}\vp(i!),&
 P&=1+D_{\max}-D_0.\label{eq:precision-all}
\end{align}
For any two supports of size $t$ whose residue multisets agree modulo
$p^P$, their optimal profiles $C_E(r)$ agree for every $r\ge r_0$.
\end{theorem}

\begin{proof}
The coefficients of $F_E$, each remainder vector $q_k$, and each
$\det q_J$ are integer polynomials in the exponents. Congruent
exponent multisets therefore give congruent determinants modulo $p^P$.
Their valuations agree whenever they are less than $P$.

The initial basis has determinant one and cost
$-D_0+rI_0$, where $I_0=t(t-1)/2$. For every $t$-element $J$,
$\sum_{j\in J}j\ge I_0$ and
$\sum_{j\in J}\vp(j!)\le D_{\max}$. Hence if
$\vp(\det q_J)\ge P$, its cost exceeds the initial cost:
\[
 \nu_J(r)-\nu_{J_0}(r)
 \ge P-D_{\max}+D_0+r(\textstyle\sum J-I_0)\ge1.
\]
Such a subset can never minimize, for any $r\ge r_0$. All possibly
minimizing costs, and thus all minimum-basis index sets, are identical
for the two supports. The cutoff at $r_0$ also applies to every deeper
disk. Theorem~\ref{thm:exact} finishes the proof.
\end{proof}

\begin{proof}[Proof of Theorem~\ref{thm:headline}]
The first part is Theorem~\ref{thm:residues-all}. A minimum of finitely
many affine functions has finitely many rational transition points.
The optimal root count is nonincreasing as the disk shrinks. At a
crossing, every basis that minimizes just to its left still minimizes
at equality, so the larger count is retained. This proves the profile
assertion, including its endpoint convention.

For computability enumerate all multisets of $t$ residues modulo
$p^P$. Realize repeated occurrences of a residue by distinct
nonnegative integers congruent to it, for example by adding successive
multiples of $p^P$. Apply the exact formula to each representative and
take the largest answer. Every support is covered. Allowing fewer
terms gives no larger optimum, because its support can be enlarged
and zero coefficients initially permitted. Proposition~\ref{prop:witness}
then imposes exactly $t$ terms, integer coefficients and irreducibility.
\end{proof}

\paragraph{Meaning of the result.}
Lenstra's existence theorem gives a computable upper bound. Theorem
\ref{thm:headline} gives a terminating procedure for the \emph{least}
such local geometric bound. It does not determine the largest possible
number of roots in a specified finite extension $K/\Qp$ throughout
$K^\times$: that is a different problem, because the roots here may
require arbitrary ramification and the coefficient optimization is
local to one disk. In particular we do not claim to solve the separate
computability question for Lenstra's global constant $B(k,K)$.

\section{A smaller uniform cutoff in the exponential region}
Now put $r_c=1/(p-1)$ and restrict to $r>r_c$. The matrix can use the
simpler power columns $v_k=\remn_{F_E}T^k$, rather than falling
factorials. These columns arise from $f(\exp(bY))$ and have the same
factorial weights
\[
 w_r(k)=rk-\vp(k!)=(r-r_c)k+\frac{\ds(k)}{p-1},
\]
where $\ds(k)$ is the base-$p$ digit sum.

For $t\ge2$ define
\begin{align*}
 D&=\max_{0\le i<t}\ds(i),&H&=D-1,&d&=\lceil\log_p t\rceil.
\end{align*}
If $H=0$ set $N_*=t-1$. If $H>0$ put
\begin{equation}\label{eq:exp-cutoff}
 M=\left\lceil\frac H{p-1}\right\rceil,\qquad
 R=d+M-1,\qquad N_* =\max\{t-1,Hp^R\}.
\end{equation}
For $t=1$ set $N_*=0$.

\begin{theorem}[Depth-independent cutoff]\label{thm:exp-cutoff}
For every support of size at most $t$ and every $r>r_c$, the last
minimum coefficient of $f(\exp(bY))$ is at most $N_*$. Thus $N_*$ is
a uniform bound for roots in $D_r$ and in the open disk
$D^\circ=\{x:\vp(x-1)>r_c\}$. For fixed $p$,
$N_*=O_p(t^2\log(2t))$.
\end{theorem}

\begin{proof}
Pad the support with unused exponents if necessary so that $|E|=t$.
Work in $A=\Z_{(p)}[T]/F_E$. Write
$M_k=\sum_e c_e e^k=L(T^k)$. Normalize by the least valuation among
$M_0,\ldots,M_{t-1}$; their Vandermonde matrix is invertible, so this
least valuation is finite. The monic recurrence implies that all
moments are integral and an early moment is a unit.

Suppose a minimizing index $k\ge t$ occurs. Comparing it with an
early unit at index $i<t$, and writing $\delta=r-r_c>0$, gives
\begin{equation}\label{eq:moment-bound}
 0\le\vp(M_k)\le\delta(i-k)+\frac{\ds(i)-\ds(k)}{p-1}.
\end{equation}
Consequently $\ds(k)\le H$. If $H=0$, this is impossible for $k>0$.
Otherwise $\ds(k)\ge1$ gives
$\vp(M_k)<H/(p-1)\le M$.

Modulo $p$, the algebra is a product of
$\Fp[T]/(T-a)^{s_a}$, with $s_a\le t\le p^d$. Hence
\[
                 T^{p^{d+1}}=T^{p^d}\pmod{pA}.
\]
In any commutative $\Z_{(p)}$-algebra,
$x\equiv y\pmod{p^a}$ with $a\ge1$ implies
$x^p\equiv y^p\pmod{p^{a+1}}$, by the binomial theorem. Repeated
powering therefore gives
\[
 T^{p^{R+1}}\equiv T^{p^R}\pmod{p^M A},
 \qquad T^{p^j}\equiv T^{p^R}\pmod{p^M A}\quad(j>R).
\]
Collapse all base-$p$ digits above position $R$ into position $R$.
The resulting integer $k'$ satisfies
\[
 T^k\equiv T^{k'}\pmod{p^M A},\qquad
 k'\le Hp^R,\qquad \ds(k')\le\ds(k).
\]
Carries only lower digit sum. If $k>N_*$, then $k'<k$ and therefore
$w_r(k')<w_r(k)$. Because $\vp(M_k)<M$, the moment congruence gives
$\vp(M_{k'})=\vp(M_k)$. This contradicts minimality.

Every polynomial has finitely many roots, so a bound on all rational
closed disks strictly inside $D^\circ$ also bounds the open disk.
Finally $D\le(p-1)d$, $M\le d$ and
$Hp^R\le(p-1)d\,p^{2d-1}$, which is $O_p(t^2\log(2t))$.
\end{proof}

\begin{theorem}[Uniform residue precision and the open endpoint]\label{thm:exp-precision}
Put
\[
 S_0=\sum_{i=0}^{t-1}\ds(i),\qquad
 P_* =\max\left\{1,\left\lceil\frac{S_0-t+1}{p-1}\right\rceil\right\}.
\]
The whole profile $C_E(r)$ for $r>r_c$, and the maximum on $D^\circ$,
depend only on the residue multiset of $E$ modulo $p^{P_*}$.
Using the power matrix $v_0,\ldots,v_{N_*}$, the open-disk answer is
the largest index in any basis minimizing, lexicographically,
\begin{equation}\label{eq:open-rule}
 \left(\vp(\det v_J)+\frac1{p-1}\sum_{j\in J}\ds(j),
                     \sum_{j\in J}j\right).
\end{equation}
\end{theorem}

\begin{proof}
The initial power basis $J_0$ has determinant one, index sum
$I_0=t(t-1)/2$, and digit sum $S_0$. Every other $t$-subset has strictly
larger index sum, while the digit sum of any $t$ distinct nonnegative
indices is at least $t-1$. For a minimizing $J\ne J_0$, comparison
with $J_0$ at $r=r_c+\delta$ gives
\[
 \vp(\det v_J)
 <\frac{S_0-\sum_{j\in J}\ds(j)}{p-1}
 \le\frac{S_0-t+1}{p-1}\le P_*.
\]
Thus determinants with valuation at least $P_*$ can be discarded;
every relevant valuation is fixed by exponent residues modulo
$p^{P_*}$, exactly as in Theorem~\ref{thm:residues-all}.

There are finitely many basis costs. For sufficiently small positive
$\delta$, minimizing their intercept plus $\delta$ times their slope
is equivalent to the lexicographic rule \eqref{eq:open-rule}. This
single final chamber computes the open-disk maximum. For the upper
bound, apply the closed-disk formula at sufficiently small depths to
include the finitely many roots of any proposed polynomial. For
attainment, take a rational dual-row witness in that chamber. It has
the stated number of roots in the open disk, and the upper bound
excludes any more. Proposition~\ref{prop:witness} applies there too.
\end{proof}

\begin{remark}[A bounded representative theorem]
Every residue multiset in Theorem~\ref{thm:exp-precision} has a
representative consisting of $t$ distinct exponents in
$[0,tp^{P_*}-1]$. Thus an open-disk extremizer over all supports
may be sought in this explicitly bounded interval, even though the
original problem imposes no degree bound. The all-depth theorem has
the analogous bound $tp^P-1$. These bounds concern existence of an
extremizing support, not the size of a given polynomial's exponents.
\end{remark}

\section{The first sharp staircase beyond the residue characteristic}
\label{sec:sharp-range}

The finite algorithm can be complemented by an explicit answer in a range
where the number of terms exceeds the residue characteristic.  Write
\[
 \mathcal D_p=\{x\in\overline{\mathbf Q}_p:
                  v_p(x-1)>1/(p-1)\}.
\]
Coefficients and roots may lie in arbitrary finite extensions of
$\mathbf Q_p$, and zeros are counted with their algebraic multiplicities.
The maximum in the following theorem is over all supports of the indicated
size, as well as over coefficients.

\begin{theorem}[A sharp staircase]\label{thm:sharp-staircase}
Let $p$ be an odd prime and let $p+2\le t\le 2p-1$.  The largest possible
number of zeros in $\mathcal D_p$ of a nonzero polynomial with at most $t$
terms is
\begin{equation}\label{eq:sharp-staircase}
 B_p(t)=p^2+\left\lfloor\frac{t-p-2}{2}\right\rfloor.
\end{equation}
The same number is attained by an integer polynomial with exactly $t$
terms which is irreducible over $\mathbf Q$.  In particular, all its zeros
are simple.  Replacing $x-1$ by $x/A-1$, for any nonzero algebraic
$p$-adic number $A$, leaves the maximum unchanged.
\end{theorem}

For example, the maxima at $p=7$ and $t=9,10,11,12,13$ are respectively
$49,49,50,50,51$.  The first value extends the repository's $p+2$-term
theorem; the successive increases require two additional terms each.
For $p=3$ the range consists only of $t=5$, so the assertion there is the
already-known value nine.  The new staircase begins at $p\ge5$.

We prove a more precise assertion about supports.  If $E$ is a finite set
of distinct nonnegative exponents, put
$s_a=\#\{n\in E:n\equiv a\pmod p\}$.

\begin{theorem}[One large residue class]\label{thm:large-residue-class}
Suppose $|E|\le2p-1$, at least two of the $s_a$ are positive, and one
residue class has multiplicity $m>p$.  Put
$q=\max\{s_a:a\text{ is not that class}\}$ and $d=\min(m-p,q)$.
The exact maximum number of zeros in $\mathcal D_p$ for nonzero
polynomials supported on $E$ is
\begin{equation}\label{eq:large-profile}
                   p^2+d-1.
\end{equation}
Every exponent of $E$ can be required to occur, and an attaining
polynomial can be chosen with rational coefficients and simple nonzero
zeros.  This formula makes no assertion about a support confined to a
single residue class.
\end{theorem}

\subsection{A normalized moment test}

We record the analytic reduction needed in both directions.  For
$f(X)=\sum_{n\in E}c_nX^n$, let $s=|E|$ and
\[
 P(T)=\prod_{n\in E}(T-n),\qquad
 A=\mathbf Z_{(p)}[T]/P(T),\qquad
 M_k=\sum_{n\in E}c_n n^k.
\]
The functional $L(T^k)=M_k$ is well-defined on $A$.  Because $P$ is
monic, every $M_k$ is an integral linear combination of
$M_0,\ldots,M_{s-1}$.  The Vandermonde matrix is invertible over
$\mathbf Q$, so these first $s$ moments cannot all vanish when $f\ne0$.
Divide them by one of least valuation.  All moments are now integral,
and at least one of the first $s$ is a unit.

For rational $r>1/(p-1)$, choose an algebraic $z$ with $v_p(z)=r$.
The substitution $X=\exp(zY)$ identifies the closed unit disk in $Y$
with $v_p(X-1)\ge r$, and preserves multiplicity.  Its coefficients are
$M_kz^k/k!$.  If $s_p(k)$ denotes the base-$p$ digit sum and
$\delta=r-1/(p-1)>0$, their valuations are
\begin{equation}\label{eq:sharp-weight}
 v_p(M_k)+w_r(k),\qquad
 w_r(k)=\delta k+\frac{s_p(k)}{p-1}.
\end{equation}
The number of geometric zeros on the closed disk is the largest index
at which these valuations take their minimum.  This is the
Weierstrass degree of the restricted power series.

Suppose $s\le2p-1$ and such a minimizing index $k$ satisfies $k\ge s$.
Every $i<s$ has $s_p(i)\le p-1$.  Comparison with an early unit moment
gives
\[
 0\le v_p(M_k)
 \le\delta(i-k)+\frac{s_p(i)-s_p(k)}{p-1}.
\]
Consequently
\begin{equation}\label{eq:sharp-small-moment}
               s_p(k)\le p-2,\qquad v_p(M_k)<1.
\end{equation}
In particular, congruences modulo $p$ suffice to exclude a large
minimizing index.  If in $A/pA$ a power $T^k$ is a linear combination
of powers $T^\ell$ with
$w_r(\ell)<w_r(k)$, lift the relation to $A$.  The error has moment
valuation at least one.  By \eqref{eq:sharp-small-moment}, at least one
participating $M_\ell$ has valuation no larger than $v_p(M_k)$; its
series coefficient then has strictly smaller valuation.  This is a
contradiction.

\subsection{The upper bound}

First suppose every residue multiplicity is at most $p$.  In
\[
 R=A/pA\simeq\prod_{a:s_a>0}\mathbf F_p[T]/(T-a)^{s_a}
\]
we have $T^{p^j}=T^p$ for every $j\ge1$.  A power of digit sum at most
$p-2$ therefore has the same image as $T^{i+pj}$ with
$i+j\le p-2$.  If its original index is at least $p^2$, this replacement
strictly decreases the index and does not increase the digit sum.
The normalized moment test excludes it.  Hence in this case there are
fewer than $p^2$ zeros in every closed disk under consideration.

Next suppose one multiplicity $m$ exceeds $p$ and another is positive.
It is unique, and $m\le2p-2$.  Multiplication of $f$ by a suitable
nonnegative power of $X$ allows its residue class to be taken as zero;
this does not change any zero in the disk.  In $R$ put
\[
                  U=T^p,\qquad Z=T^{p^2}.
\]
On every factor other than the zero factor, $U=Z=a$; on the zero factor,
$U=T^p$, $Z=0$, and $U^2=0$.  Every higher Frobenius power equals $Z$.
Thus a candidate of digit sum $h\le p-2$ has image
$T^iU^jZ^\ell$, where $i+j+\ell=h$.

If $\ell=0$, its compressed index is below $p^2$.  If
$\ell\ge1$ and $j+\ell\ge2$, its image equals
$T^iU^{j+\ell}$: this follows factor by factor, using $U^2=Z=0$
on the zero factor.  Again the replacement index is below $p^2$ and
has the same digit sum.  The only remaining candidate above $p^2-1$
is therefore
\begin{equation}\label{eq:sharp-exceptional-index}
                        k=p^2+h-1.
\end{equation}

If $h>m-p$, then $T^{h-1}Z=T^{p+h-1}$ in $R$, since both vanish on
the zero factor.  This is a lower-index replacement.  If $h>q$, then
$Z(T-Z)^{h-1}=0$ in $R$: on every nonzero factor the nilpotent
$T-a$ has order at most $q\le h-1$.  Expanding gives
\[
 T^{h-1}Z
   =\sum_{j=2}^{h}(-1)^j\binom{h-1}{j-1}T^{h-j}Z^j
   =\sum_{j=2}^{h}(-1)^j\binom{h-1}{j-1}T^{h-j}U^j.
\]
Every term on the right has digit sum $h$ and index below $p^2$.
The moment test excludes \eqref{eq:sharp-exceptional-index} unless
$h\le\min(m-p,q)$.  This proves the upper bound in
Theorem~\ref{thm:large-residue-class}.

For the unrestricted theorem a support confined to one residue class
causes no difficulty.  Factor out its smallest monomial and divide all
exponent differences by their largest common power $p^a$.  The
polynomial becomes $X^b g(X^{p^a})$, where $g$ occupies at least two
residue classes.  On $\mathcal D_p$ the power map is injective, because
the logarithm identifies it with multiplication by $p^a$, and it
preserves multiplicities.  Its image is contained in $\mathcal D_p$.
Thus the bound for $g$ is also a bound for the original polynomial.
The case of a single monomial has no zeros in the disk.

For any remaining support of size at most $t$, either the bound is
less than $p^2$, or
$2\min(m-p,q)\le(m-p)+q\le t-p$.  Therefore
\[
 p^2+\min(m-p,q)-1
 \le p^2+\left\lfloor\frac{t-p}{2}\right\rfloor-1,
\]
which is \eqref{eq:sharp-staircase}.  Every polynomial has finitely many
roots, so bounds on every rational closed disk of depth greater than
$1/(p-1)$ also bound the roots in the open disk $\mathcal D_p$.

\subsection{A functional which attains the bound}

Let $1\le d\le p-2$ and suppose that the support has at least $p+d$
exponents in residue zero and at least $d$ in a residue
$a\in\mathbf F_p^\times$.  Project its residue algebra onto
\[
 R'=\mathbf F_p[T]/T^{p+d}\ \times\
                    \mathbf F_p[T]/(T-a)^d.
\]
For $F=(F_0,F_a)$ in this quotient define
\begin{equation}\label{eq:sharp-functional}
 \begin{split}
 \lambda(F)={}&[Y^{d-1}]F_a(a+Y)
       -[T^{d-1}]F_0(T)-da[T^d]F_0(T)\\
       &\hspace{32mm}-a[T^{p+d-1}]F_0(T).
 \end{split}
\end{equation}
All indicated coefficients are well-defined in the truncated factors.
Set $C=p^2+d-1$.  Lucas's congruence gives
\begin{equation}\label{eq:sharp-functional-values}
 \begin{aligned}
 \lambda(T^k)&=0&&\text{if }s_p(k)<d,\\
 \lambda(T^k)&=0&&\text{if }s_p(k)=d\text{ and }k<C,\\
 \lambda(T^C)&=a\ne0.&
 \end{aligned}
\end{equation}
Here is a direct check.  The first term of
\eqref{eq:sharp-functional} on $T^k$ is
$\binom{k}{d-1}a^{k-d+1}$.  Since $d-1<p$, its binomial factor
vanishes unless the units digit of $k$ is at least $d-1$.
With digit sum less than $d$, the only exception is $k=d-1$,
cancelled by the second term.  With digit sum $d$ and $k<C$, the
only exceptions are $k=d$ and $k=p+d-1$, cancelled by the third and
fourth terms.  At $C$ the first term equals $a^{p^2}=a$, and the
other three terms vanish.  This verification also covers $d=1$.

Compose $\lambda$ with the residue projection and lift its values on
$1,T,\ldots,T^{s-1}$ to ordinary integers.  This defines an integral
functional $L$ on $A$, and hence rational coefficients $c_n$ by
Vandermonde inversion.  Its moments reduce to
\eqref{eq:sharp-functional-values}.  Choose, explicitly,
\[
 r=\frac1{p-1}+\frac1{2(p-1)C}.
\]
Then
$w_r(C)=(d+\tfrac12)/(p-1)<\min(1,(d+1)/(p-1))$.
All moments of smaller digit sum, and all preceding indices of digit
sum $d$, are divisible by $p$ and therefore have coefficient valuation
at least one.  Indices of greater digit sum have weight at least
$(d+1)/(p-1)$; later indices of the same digit sum have strictly
larger weight.  The $C$-th moment is a unit.  Thus $C$ is the unique
minimizing coefficient index, giving exactly $C$ zeros on this closed
disk, and therefore at least $C$ in $\mathcal D_p$.

For Theorem~\ref{thm:large-residue-class}, take
$d=\min(m-p,q)$ and choose a residue attaining $q$.  The quotient
above is available.  For Theorem~\ref{thm:sharp-staircase}, put
$d=\lfloor(t-p)/2\rfloor$ and use the support
\begin{equation}\label{eq:sharp-support}
 E_{p,d}=\{jp:0\le j<p+d\}\ \cup\
                   \{1+jp:0\le j<d\}.
\end{equation}
It has $p+2d$ exponents.  If this is $t-1$, add any new nonnegative
exponent.  The previously obtained root count survives a sufficiently
small nonzero coefficient on that exponent.

For completeness, both the full-support and simplicity requirements
can be enforced without relying on genericity.  First clear
denominators.  Perturb each permitted integer coefficient by a
multiple of a sufficiently large power $p^H$.  Each moment changes
by a number of valuation at least $H$, uniformly in its index, so
\eqref{eq:sharp-weight} shows that the strict minimizing index is
preserved.  Require all permitted coefficients to be nonzero.
After removing the smallest monomial, choose a prime $\ell\ne p$.
The Chinese remainder theorem additionally makes the leading
coefficient prime to $\ell$, all other coefficients divisible by
$\ell$, and the constant coefficient divisible by $\ell$ but not by
$\ell^2$.  The resulting polynomial is Eisenstein at $\ell$ and hence
irreducible over $\mathbf Q$.  In characteristic zero its roots are
simple.  When the original support contains a nonzero smallest
exponent, irreducibility is asserted after that monomial is removed;
the supports \eqref{eq:sharp-support} contain zero.  Finally, scaling
$X$ by $A$ proves the statement for an arbitrary center.

\subsection{Independent finite checks}

The script \texttt{code/verify\_sharp\_range.py} implements the residue
algebra directly by jets $\binom{k}{j}a^{k-j}$, using only arithmetic
in $\mathbf F_p$.  It scans powers by digit sum and then by index;
this is independent of the rational determinant implementation used
for the general algorithm.  It exhausts all residue-multiplicity
vectors at $p=5$ for $t=7,8,9$, checks larger-prime examples and
deterministically seeded profiles, and verifies
\eqref{eq:sharp-functional-values} separately.  For $p=5,t=9$ it also
constructs an integer Eisenstein witness, expands it directly at
$X=1$, and verifies that its Newton weights have last minimum at
$26$ at depth $53/208$.  These checks illustrate and test the argument;
the proofs above apply to every prime and every support in the stated
range.

\section{Exact binary profiles beyond the previous threshold}\label{sec:binary-profile}
For $t\ge1$ and $r>1$, write $B_2(t,r)$ for the maximum number of
geometric roots in $v_2(X-1)\ge r$ of a nonzero polynomial with at most
$t$ nonzero terms. Coefficients may lie in $\mathbb C_2$; allowing only
rational coefficients gives the same maxima below. Roots are counted with
multiplicity for the upper bounds, and all lower bounds have distinct roots.
Write $B_2^\circ(t)$ for the corresponding maximum in $v_2(X-1)>1$.

\begin{theorem}[Five, six, and seven terms over the binary disk]\label{thm:binary-profile}
For every real $r>1$,
\[
 B_2(5,r)=4,
 \qquad
 B_2(6,r)=
 \begin{cases}8,&1<r\le4/3,\\5,&r>4/3,\end{cases}
 \qquad
 B_2(7,r)=
 \begin{cases}8,&1<r\le3/2,\\6,&r>3/2.\end{cases}
\]
Consequently
\[
 B_2^\circ(5)=4,\qquad B_2^\circ(6)=B_2^\circ(7)=8.
\]
Each displayed maximum is attained by an integer polynomial with exactly
$t$ terms and with all of its roots in the indicated disk simple.
\end{theorem}

\begin{proof}
The cutoff theorem has $D=2$, $H=1$, $M=1$, $d=3$ and $N=8$ for
$t=5,6,7$. After normalizing the first $t$ moments to be integral with
at least one unit, any minimizing index $k\ge t$ has
$s_2(k)\le1$ and $v_2(M_k)<1$. Thus the only such index is $k=8$.
The coefficients are weighted by
\[
 w_r(k)=rk-v_2(k!)=(r-1)k+s_2(k).
\]
For $t=6$, the largest of the early weights is $w_r(5)=5r-3$ for
$r>1$. Since $w_r(8)=8r-7$, the index $8$ is impossible if $r>4/3$.
For $t=7$, the largest early weight is $w_r(6)=6r-4$, so the index $8$
is impossible if $r>3/2$. This proves the upper bounds of $8$ and the
stronger bounds beyond those respective endpoints.

For $t=5$, we exclude the index $8$ at every $r>1$. In the exponent
algebra modulo two, write
\[
 \overline A=\mathbb F_2[T]/\bigl(T^{a}(T-1)^{5-a}\bigr).
\]
If $1\le a\le4$, then $T^8=T^4$ in $\overline A$: on either primary
factor the fourth and eighth powers are the corresponding scalar zero
or one. Hence $M_8\equiv M_4\pmod2$. Because a minimizing $M_8$ would
have valuation below one, the two moments would have equal valuation.
But $w_r(8)-w_r(4)=4(r-1)>0$, a contradiction.
If $a=5$, then $T^8=0$, making $v_2(M_8)\ge1$, again impossible.
If $a=0$, then $T^8=1$, and comparison with $M_0$ gives the same
contradiction because $w_r(8)>0$.
This proves the upper bound four. Allowing fewer than five terms cannot
increase it: one may pad a support by unused columns before applying
the argument. The same padding convention covers the earlier bounds.

For sharpness in the two branches with value eight, use
\begin{align*}
 F_6(X)&=5X^8-28X^6+70X^4-140X^2+128X+65501,\\
 F_7(X)&=3X^8-28X^6+210X^4-448X^3+420X^2-192X+65571.
\end{align*}
Their exact expansions at one are
\begin{align*}
 F_6(1+Y)&=2^{16}+112Y^5+112Y^6+40Y^7+5Y^8,\\
 F_7(1+Y)&=2^{16}+56Y^6+24Y^7+3Y^8.
\end{align*}
The lower Newton polygon of the first has vertices $(0,16),(5,4),(8,0)$.
It gives five roots of valuation $12/5$ and three of valuation $4/3$.
The second has vertices $(0,16),(6,3),(8,0)$, giving six roots of
valuation $13/6$ and two of valuation $3/2$.
Thus both endpoint inequalities retain all eight roots.
These roots are simple. Indeed the following explicit identities hold:
\[
 \begin{split}
 &(X^6+X^3-1)F_6+
 (X^7+X^5+X^4+X^3+X^2+X)F_6'\equiv1\pmod3,\\
 &X^5F_7+(-2X^6-2X^4+X^2+2)F_7'\equiv1\pmod5.
 \end{split}
\]
In each identity the leading coefficient of $F_i$ is a unit modulo the
indicated prime, so coprimality after reduction implies squarefreeness
over $\mathbb Q$.

For the remaining branches, take $(X-1)^{t-1}-2^K$ with $K$ a sufficiently
large positive integer, specifically $K/(t-1)\ge r$. This polynomial has
exactly $t$ nonzero terms and $t-1$ distinct roots, each of depth
$K/(t-1)$. For the open disk choose a strict inequality. All assertions
for real $r$ also follow directly from these witnesses and from the
weight comparisons, or by rational approximation in the cutoff theorem.
\end{proof}

The five-term value stays at four, while the sixth permitted term doubles
the open-disk maximum to eight. The seventh term then increases the
closed depth through which eight roots can be retained, from $4/3$ to
$3/2$, without increasing the open-disk maximum. These are geometric
root counts in $\mathbb C_2$, not counts of roots lying in $\mathbb Q_2$.
The earlier repository result through $t=p+2$ only reaches four terms
at $p=2$; the present five-, six-, and seven-term formulas lie beyond
that range.


\section{Exact calculations and reproducible artifacts}
The supplied implementation uses Python's integer and rational
arithmetic. Its two realizations of the theorem are:
\begin{itemize}
\item falling-factorial columns and Lenstra's cutoff, valid at every
positive rational depth;
\item power columns and Theorem~\ref{thm:exp-cutoff}, valid throughout
the exponential region, including its open limiting rule.
\end{itemize}
It returns a minimizing basis, the indices present in all minimizing
bases collectively, and a rational-coefficient witness, cleared to
integers. For modest polynomial degrees an independent routine expands
$f(1+Y)$ directly and checks its Newton polygon. No numerical root
approximation or probabilistic primality test is used.

\subsection{A nine-term profile at five}
For the support
\[
 E=\{0,1,5,6,10,15,20,25,30\}
\]
the exact profile obtained by the basis inequalities is
\begin{center}
\begin{tabular}{cc}
\toprule
Depth $r$ & $C_E(r)$\\
\midrule
$1/4<r\le5/19$ & 26\\
$5/19<r\le4/15$ & 25\\
$4/15<r\le1/2$ & 10\\
$r>1/2$ & 8\\
\bottomrule
\end{tabular}
\end{center}
Every endpoint is handled by tied minimum bases, not by floating-point
rounding. The accompanying JSON file records the complete computation.
The explicit sharp-range witness is checked separately at
$r=53/208>1/4$, where it has exactly 26 roots. Its coefficients are
Eisenstein at two, providing a short independent certificate of
irreducibility and simplicity.

\subsection{What the verification establishes}
The archive contains separate checkers for the general matrix rule,
the odd-prime staircase, and the binary profiles. The matrix checker
compares the quotient-algebra result with exhaustive determinants of
the full, original binomial coefficient matrix for small supports.
The staircase checker enumerates residue multiplicities and tests the
explicit residue functional, while the witnesses are checked from
ordinary polynomial expansions. The torsion companion has its own
exact checks. The result files give the actual case counts and inputs;
\texttt{notes/VERIFICATION.txt} summarizes them.

These finite checks test the implementations and the displayed finite
certificates. The arbitrary-prime and arbitrary-support assertions
rest on the proofs in this paper. No Lean or other proof-assistant
verification of these new theorems is claimed.

\section{What this resolves and what remains open}
The all-depth determinant theorem supplies an exact algorithm for the
optimal geometric local fewnomial constant. The Frobenius theorem
gives one finite classification valid at every depth in the exponential
region, resolving the arbitrary-term extension asked for by the recent
repository reports. The sharp staircase and binary theorems provide
closed formulas beyond the previous $p+2$-term threshold.

Several substantial questions remain. The precision and degree bounds
used for finite enumeration are far from optimized. A closed formula
for the open-disk optimum at all $p,t$ is not supplied, and constraints
requiring roots to lie in a fixed finite extension of $\Qp$ are not
recovered from geometric root counts.
The related singular-modulus six-term problem is also not decided by
these local bounds alone; it requires information about where its
conjugates lie.

We also investigated the complete-bipartite canonical-line question
for Ehrhart polynomials \cite{HKM,Kolbl}, following the repository's
finite-width proof. We did not obtain the missing all-width positivity
argument. The accompanying research-status note records that attempt
and its precise obstacle. Neither this paper nor the torsion companion
claims a proof of a prominent named conjecture. Their new conclusions
are the explicitly stated local classification and growth theorems.

\begin{thebibliography}{99}\raggedright
\bibitem{Lenstra} H. W. Lenstra, Jr., \emph{On the factorization of
lacunary polynomials}, in Number Theory in Progress, Vol.~1,
de Gruyter, 1999, pp.~277--292.
\href{https://pub.math.leidenuniv.nl/~lenstrahw/PUBLICATIONS/1999b/art.pdf}{Author's scan}.
\bibitem{PoortenRumely} A. J. van der Poorten and R. S. Rumely,
\emph{Zeros of $p$-adic exponential polynomials II}, J. London Math.
Soc. (2) 36 (1987), 1--15.
\href{https://doi.org/10.1112/jlms/s2-36.1.1}{doi:10.1112/jlms/s2-36.1.1}.
\bibitem{AvendanoKrick} M. Avenda\~no and T. Krick,
\emph{Sharp bounds for the number of roots of univariate fewnomials},
J. Number Theory 131 (2011), 1209--1228.
\href{https://arxiv.org/abs/1008.4808}{arXiv:1008.4808}.
\bibitem{Battistella} L. Battistella, K. K\"uhn, A. Kuhrs, M. Ulirsch
and A. Vargas, \emph{Buildings, valuated matroids, and tropical linear
spaces}, J. London Math. Soc. 109 (2024), no.~1, e12850.
\href{https://doi.org/10.1112/jlms.12850}{doi:10.1112/jlms.12850}.
\bibitem{ProveIt} ProveIt contributors, \emph{Beyond the Riemann
hypothesis for Dirichlet $L$-functions: consequences, transfers and
limits}, repository snapshot cited in Section 1.
\href{\repo/Article/README.md}{Article inventory};
\href{\repo/Article/sections/smd.tex}{singular-modulus and local root section};
\href{\repo/Article/sections/cost.tex}{torsion cost section}.
\bibitem{RankPacket} ProveIt research packet, \emph{An exact rank rule
for sparse roots in closed $p$-adic disks}, 27 November 2026.
\href{\repo/Research/padic-support-rank/report.tex}{Pinned source}.
\bibitem{SharpPacket} ProveIt research packet, \emph{Sharp root counts
for sparse polynomials in the open $p$-adic exponential disk},
27 November 2026.
\href{\repo/Research/sparse-padic-roots-and-graph-towers/fewnomial/paper.tex}{Pinned source}.
\bibitem{Sernau} L. Sernau and collaborators, \emph{GRH for Dirichlet
$L$-functions, in Lean 4}, commit
\texttt{72fead102f663d5c7a407964a13dc0cf9d977207}.
\href{https://github.com/lsernau/grh-lean}{Repository}.
\bibitem{Sutherland} A. V. Sutherland, \emph{Independent verification
of the shorter GRH proof}, inspected 27 November 2026.
\href{https://math.mit.edu/~drew/shortgrh/}{Verification report}.
\bibitem{HKM} A. Higashitani, M. Kummer and M. Micha\l ek,
\emph{Interlacing Ehrhart polynomials of reflexive polytopes},
Selecta Math. 23 (2017), 2977--2998.
\href{https://arxiv.org/abs/1612.07538}{arXiv:1612.07538}.
\bibitem{Kolbl} M. K\"olbl, \emph{On a conjecture concerning the roots
of Ehrhart polynomials of symmetric edge polytopes from complete
multipartite graphs}, 2024.
\href{https://arxiv.org/abs/2404.02136}{arXiv:2404.02136}.
\end{thebibliography}
\end{document}
